\paragraph{Quadratic rational first-integral pencils.}
A bounded list theorem also applies when the denominator of a first
integral depends on the candidate value.  The square-value input below
is the classical function-field theorem of Pasten and Wang
\cite[Theorem~3]{pasten-wang-buchi}; the argument keeps track of the
height needed to exclude its positive-characteristic exceptions.

\begin{theorem}\label{thm:quadratic-rational-pencil-list}
Let $\mathbb K$ be a field of characteristic $p>\max\{2,10D\}$,
where $0\le D<n$, and let $N,T\in\mathbb K(X)[u]$ be coprime with
$\max\{\deg_u N,\deg_u T\}=2$.
Consider the polynomial sections
\[
 \mathcal B=\{P\in\mathbb K[X]:\deg P\le D,
       \ N(X,P)=cT(X,P)\text{ for some }c\in\mathbb K\}.
\]
On any $n$ distinct evaluation coordinates in $\mathbb K$, every received word has
at most
\[
 \max\{40,\ \lfloor 1/\eta\rfloor\}
\]
members of $\mathcal B$ agreeing on at least $D+\eta n$ coordinates,
for every $\eta>0$.  The received symbols may lie in an extension
field.  This is a bound for the indicated section family, not for
all Reed--Solomon codewords.
\end{theorem}
\begin{proof}
Extend constants to an algebraic closure $k$ and put $K=k(X)$.
A section has a unique label: otherwise $N(P)=T(P)=0$, contrary
to coprimality.  Each label has at most two sections.  Thus a
selected list with fewer than twenty-one labels has size at most
forty.  Suppose henceforth that twenty-one labels occur.

\emph{Reconstructing a low-height pencil.}
Choose five sections $P_i$ with distinct labels $c_i$.  They are
distinct elements of $K$.  The matrix with rows
\[
 (P_i^2,P_i,1,-c_iP_i^2,-c_iP_i,-c_i)
\]
has rank five over $K$.  Indeed, another kernel vector defines
$(\widetilde N,\widetilde T)$ of degrees at most two whose
cross-difference $\widetilde NT-N\widetilde T$ has degree at most
four and vanishes at all five $P_i$.  It is zero; coprimality and
the exact degree-two maximum make the pairs proportional.
The signed maximal minors therefore give a polynomial coefficient
representative $(n_2,n_1,n_0,t_2,t_1,t_0)$ for $(N,T)$ with degrees
bounded respectively by
\[
 (4D,5D,6D,4D,5D,6D).
\]
Replace $(N,T)$ by this common polynomial representative.
Consequently the coefficients of the discriminant
\[
 F(C)=(n_1-Ct_1)^2-4(n_2-Ct_2)(n_0-Ct_0)
\]
have $X$-degree at most $10D$.  The polynomial $F$ has degree at
most two in $C$ and is not a square in $K(C)$: the coprime
rational map $N/T$ has degree two, so $N-CT$ is irreducible over
$K(C)$ as a polynomial in $u$.

Every section gives
\[
 F(c_i)=\bigl(2(n_2-c_it_2)P_i+n_1-c_it_1\bigr)^2.
\]
This remains valid for a linear specialized fiber.  Choose $c_0$
with $F(c_0)=s^2\ne0$, possible since $F$ has at most two zeros,
and form
\[
 H(Z)=\frac{Z^2F(c_0+1/Z)}{F(c_0)}.
\]
It is monic quadratic.  Its projective coefficient vector is
$[F(c_0):F'(c_0):[C^2]F]$, so its function-field height is at most
$10D<p$.  At the twenty distinct constants
$Z_i=(c_i-c_0)^{-1}$ it takes square values.

\emph{Applying the square-value theorem.}
For a monic quadratic over the genus-zero field $K$, the
Pasten--Wang theorem has threshold $M>11\cdot2-3=19$.
After removing constant polynomial factors, it requires every
remaining irreducible factor to be separable and outside $K^p[Z]$;
with twenty square values it makes every such factor have
multiplicity two.  These hypotheses hold here.  Degree at most
two and $p>2$ give separability.  Monic factor heights add, by
multiplicativity of the local Gauss valuation, so every factor
has height below $p$.  A monic factor in $K^p[Z]$ with nonconstant coefficients
would have a nonconstant coefficient which is a $p$th power,
with height at least $p$, a contradiction.
If $H$ has nonconstant coefficients, the multiplicity conclusion
and its total degree two force it to be a square; a positive-degree
constant factor cannot remain in that case.  But the reciprocal
substitution and its square multiplier would then make $F$ a
square in $K(C)$.  Therefore $H\in k[Z]$, and undoing the
substitution yields
\[
 F(C)=s^2q(C),\qquad q\in k[C],\quad \deg q\le2.
\]
The polynomial $q$ is nonsquare, hence squarefree of degree one or
two over the algebraically closed field $k$.

\emph{A constant conic and its sections.}
Write $N-CT=a(C)u^2+b(C)u+d(C)$.  In $K(u)$ put
$c=N(u)/T(u)$ and $v=(2a(c)u+b(c))/s$.  Then
\[
 v^2=q(c),\qquad K(u)=K(c,v).
\]
The smooth projective constant conic $v^2=q(c)$ is rational:
parametrizing by lines through a constant point gives its function
field $k(t)$.  Hence $K(u)=K(t)$, so $u$ is a degree-one rational
function of $t$,
\[
 u=\frac{A(X)t+B(X)}{C(X)t+E(X)},\qquad AE-BC\ne0.
\]
Every polynomial section corresponds to a constant parameter in
$\mathbb P^1(k)$.  Indeed its $c$ is constant and its $v$ has square
$q(c)\in k$, hence is constant too.  The isomorphism of smooth
projective curves handles degree-drop fibers and parameters at
infinity.  Distinct sections give distinct parameters.  Polynomial
members of this M\"obius family need not form an affine line.

\emph{Local collisions and agreement.}
At each evaluation coordinate $x$, scale the $2\times2$ coefficient
matrix over the discrete valuation ring at $X=x$ to make its
entries integral with at least one unit.  Its reduction has rank
one or two.  Rank two gives distinct values at all selected
parameters.  In rank one, all parameters except the unique kernel
point have the same projective image.  That image is finite since
at least one nonkernel selected section is a polynomial.  Thus
all but at most one selected value coincide.

Let $L$ be the selected list size, let $m_x$ be the number of
selected candidates matching the received symbol at $x$, and let
$b_x$ count all agreeing candidate pairs there.  The local description
gives $m_x\in\{0,1,L-1,L\}$, and in every case
\[
 (m_x-1)(L-1)\le 2b_x.
\]
For $m_x=L-1$ this follows from $b_x\ge\binom{L-1}{2}$;
for $m_x=L$ use $b_x=\binom L2$; the other cases are immediate.
Every distinct polynomial pair agrees at at most $D$ coordinates.
Summing, and writing $a$ for the minimum agreement count, gives
\[
 (La-n)(L-1)\le 2\sum_x b_x
       \le D L(L-1),\qquad L(a-D)\le n.
\]
Thus $a-D\ge\eta n$ implies $L\le1/\eta$.
Together with the forty-section
case, this proves the stated bound.  Extending the constants did
not change any agreement count or polynomial degree, so the
conclusion holds over the original field.
\end{proof}
